\section{Discussion and Limitations}
\subsection{When separation helps, and when it does not}
The proposed boundary is most promising when data coordination or input waits
are large enough to leave useful execution opportunities unfilled, but task
descriptors are cheap enough to queue. It can be less effective for long,
purely computational tasks; for short graphs dominated by startup; or for
workloads whose intermediates overwhelm available storage. A larger window
cannot manufacture bandwidth once senders, memory channels, or disks saturate.
An eager preparation policy may outperform deferred preparation when prefetch
latency matters more than prepared-data occupancy.

The logical scheduler currently uses a least-loaded choice rather than a
comprehensive locality cost model. Moving replica choice out of placement
simplifies one boundary but may incur more remote reads. Integrating locality
hints would require stating whether they are advisory, how stale hints are
handled, and whether they reintroduce synchronous source coordination. The
present architecture permits that extension; the present measurements do not
prove that it is unnecessary.

\subsection{Recovery costs and storage limits}
Worker-local intermediates expose a different recovery tradeoff from
background backup. Requested-output durability does not automatically protect
all intermediate state. Likewise, fsync on a controller-local filesystem
provides a host-specific contract, not replicated durability across host loss.
A future controller replication design would need explicit authority transfer,
log ordering, and failure fencing. Adding a second controller process without
those rules would not establish the stronger guarantee.

Open workflows can retain data for consumers not yet declared. Storage
credits, spill policies, and cancellation cleanup therefore deserve separate
capacity experiments. Memory sampling and declared-memory admission do not
bound every user allocation, and the current feedback ignores some bottlenecks.
An extension using transfer-byte credits or independent compute and transfer
windows should be evaluated against the current simpler policy, including
cases where extra state and control work make it slower.

\subsection{What the new evidence establishes}
The complete ATLAS analysis closes the gap between isolated numerical kernels
and an end-to-end workflow over real scientific data. It does not establish
production collaboration readiness, heterogeneous deployment, or distributed
dataset ingestion. The fixed/elastic comparison also changes the interpretation:
the execution boundaries and their ownership are the evaluated design, while
feedback remains an implementation choice whose incremental benefit is limited
on the tested graphs.

Five-block measurements across eight distinct physical hosts provide a larger
and more controlled scaling study than counting opportunistically colocated
workers. They still lack exclusive nodes, 16+ hosts, and a controller-capacity
sweep at fixed application work. The data-heavy results caution against assuming
that independent data ownership removes every centralized or placement cost.
Cold-library sensitivity similarly makes dependency preparation part of the
deployment contract, not an incidental detail that can be omitted from a claim.

Native grouping is now exercised with compatible start-credit semantics and
expanded windows. A prepared-node assignment policy remains an unimplemented
baseline; eager preparation alone is narrower than that policy. Exact sender
traffic and physical prepared-byte accounting remain necessary for stronger
claims about network or storage savings. Worker-lifetime CPU and ordered local
traces strengthen attribution but do not replace those measurements.
Together, these results support explicit ownership of logical scheduling and
data-lifecycle service as the paper's architectural contribution.
